% =====================================================================
\chapter{Sequências e séries}

Somar infinitos números pode dar um resultado finito: $\frac12+\frac14+\frac18+\dots=1$.
Mas também pode não dar: $1+\frac12+\frac13+\dots$ cresce sem limite, embora os termos
tendam a zero. Este capítulo organiza essa questão. Primeiro estudamos \emph{sequências}
(listas infinitas de números) e seus limites; depois, \emph{séries} (somas infinitas), definidas
como limites de somas parciais, e uma coleção de testes para decidir se convergem. As séries de
potências do próximo capítulo dependem de tudo isso.

% ---------------------------------------------------------------------
\section{Sequências}

Uma \textbf{sequência}\index{sequência} é uma lista ordenada infinita $a_1,a_2,a_3,\dots$,
indicada por $\{a_n\}$. Pode ser dada por uma fórmula explícita ($a_n=\frac n{n+1}$) ou por
uma \emph{recorrência} ($a_1=1$, $a_{n+1}=\frac{a_n}2+1$).

\begin{definicao}
A sequência $\{a_n\}$ \textbf{converge} para $L$, e escrevemos $\lim_{n\to\infty}a_n=L$, se para
todo $\varepsilon>0$ existe $N$ tal que $|a_n-L|<\varepsilon$ sempre que $n>N$. Se o limite não
existe (ou é infinito), a sequência \textbf{diverge}.
\index{sequência!convergente}
\end{definicao}

\textbf{Ferramentas de cálculo.} Se $a_n=f(n)$ e $\lim_{x\to\infty}f(x)=L$, então
$a_n\to L$: tudo o que sabemos de limites no infinito (inclusive L'Hôpital, aplicado a $f$)
vale para sequências. As regras de soma, produto e quociente continuam valendo, assim como o
Teorema do Confronto. Se $a_n\to L$ e $g$ é contínua em $L$, então $g(a_n)\to g(L)$.

\begin{exemplo}{2.1}{Limites de sequências}
(a) $\dfrac{3n^2-1}{2n^2+n}\to\dfrac32$ (divida por $n^2$).

(b) $\dfrac{\ln n}{n}\to0$: por L'Hôpital em $f(x)=\frac{\ln x}x$, $\lim\frac{1/x}1=0$.

(c) $\Big(1+\dfrac3n\Big)^n\to e^3$ (o limite fundamental do capítulo~3 do Cálculo~I).

(d) $\dfrac{\cos n}{n}\to0$ pelo confronto: $-\frac1n\le\frac{\cos n}n\le\frac1n$.

(e) $r^n\to0$ se $|r|<1$; $r^n\to1$ se $r=1$; e $\{r^n\}$ diverge se $r=-1$ ou $|r|>1$.
\end{exemplo}

\textbf{Sequências monótonas.} $\{a_n\}$ é \emph{crescente} se $a_{n+1}\ge a_n$ e
\emph{decrescente} se $a_{n+1}\le a_n$; é \emph{limitada} se existe $M$ com $|a_n|\le M$ para
todo $n$. Toda sequência convergente é limitada, mas a recíproca é falsa ($(-1)^n$).

\begin{teorema}[Teorema da convergência monótona]{2.1}
\index{Teorema!da convergência monótona}Toda sequência monótona e limitada converge.
\end{teorema}

\begin{exemplo}{2.2}{Uma recorrência}
Seja $a_1=1$, $a_{n+1}=\frac{a_n}2+1$. Por indução, $a_n<2$ e $a_{n+1}-a_n=1-\frac{a_n}2>0$:
a sequência é crescente e limitada, logo converge para algum $L$. Passando ao limite na
recorrência, $L=\frac L2+1$, isto é, $L=2$.
\end{exemplo}

% ---------------------------------------------------------------------
\section{Séries infinitas}

\begin{definicao}
Dada a sequência $\{a_n\}$, a \textbf{série}\index{série} $\sum_{n=1}^\infty a_n$ tem
\textbf{somas parciais}\index{soma parcial} $S_k=a_1+a_2+\dots+a_k$. Se $S_k\to S$, a série
\textbf{converge} e tem \emph{soma} $S$; caso contrário, \textbf{diverge}.
\end{definicao}

\begin{teorema}[Série geométrica]{2.2}
\index{série!geométrica}A série $\displaystyle\sum_{n=1}^\infty ar^{n-1}=a+ar+ar^2+\cdots$
($a\neq0$) converge se, e somente se, $|r|<1$, e nesse caso sua soma é $\dfrac a{1-r}$.
\end{teorema}
\emph{Prova.} $S_k-rS_k=a-ar^k$, logo $S_k=\frac{a(1-r^k)}{1-r}$ para $r\neq1$, e $r^k\to0$
exatamente quando $|r|<1$. $\square$

\begin{exemplo}{2.3}{Geométricas e telescópicas}
(a) $\sum_{n=1}^\infty\frac{3}{4^n}=\frac{3/4}{1-1/4}=1$.

(b) A dízima $0{,}\overline{27}=\frac{27}{100}+\frac{27}{100^2}+\dots=\frac{27/100}{1-1/100}=\frac{27}{99}=\frac3{11}$.

(c) \emph{Telescópica:}\index{série!telescópica} $\frac1{n(n+1)}=\frac1n-\frac1{n+1}$, logo
$S_k=\big(1-\tfrac12\big)+\big(\tfrac12-\tfrac13\big)+\dots+\big(\tfrac1k-\tfrac1{k+1}\big)=1-\tfrac1{k+1}\to1$.
\end{exemplo}

\textbf{A série harmônica}\index{série!harmônica} $\sum\frac1n$ diverge, embora
$\frac1n\to0$. Agrupando os termos:
\[
1+\tfrac12+\big(\tfrac13+\tfrac14\big)+\big(\tfrac15+\dots+\tfrac18\big)+\cdots
\ \ge\ 1+\tfrac12+\tfrac12+\tfrac12+\cdots,
\]
então $S_{2^k}\ge1+\frac k2\to\infty$. O crescimento é lentíssimo: são precisos mais de $12\,000$
termos para passar de $10$.

\figtikz{\begin{tikzpicture}
\begin{groupplot}[group style={group size=2 by 1,horizontal sep=1.3cm},painel,width=6.4cm,height=4.6cm,xmin=0,xmax=21,xtick={5,10,15,20},xlabel={$k$}]
\nextgroupplot[ymin=0,ymax=1.25,ytick={0.5,1},title={$\sum\frac1{2^n}$: $S_k\to1$}]
\addplot[cinza,dashed,domain=0:21]{1};
\addplot[only marks,mark=*,mark size=1.3pt,azul,samples at={1,...,20}]{1-0.5^x};
\nextgroupplot[ymin=0,ymax=4,ytick={1,2,3},title={$\sum\frac1n$: $S_k\to\infty$}]
\addplot[only marks,mark=*,mark size=1.3pt,laranja] coordinates {(1,1)(2,1.5)(3,1.8333)(4,2.0833)(5,2.2833)(6,2.45)(7,2.5929)(8,2.7179)(9,2.829)(10,2.929)(11,3.0199)(12,3.1032)(13,3.1801)(14,3.2516)(15,3.3182)(16,3.3807)(17,3.4396)(18,3.4951)(19,3.5477)(20,3.5977)};
\addplot[cinza,dashed,domain=1:21]{ln(x)+0.5772};
\end{groupplot}\end{tikzpicture}}{Somas parciais de duas séries com termos tendendo a zero. A
geométrica converge; a harmônica cresce como $\ln k$ (tracejada) e diverge.}

\textbf{Álgebra de séries.} Se $\sum a_n$ e $\sum b_n$ convergem, então $\sum(a_n\pm b_n)$ e
$\sum c\,a_n$ convergem, com as somas esperadas. Alterar ou remover um número finito de termos
não muda a convergência (só a soma).

% ---------------------------------------------------------------------
\section{Teste da divergência e teste da integral}

\begin{teorema}[Teste da divergência]{2.3}
\index{teste!da divergência}Se $\lim a_n\neq0$ (ou não existe), então $\sum a_n$ diverge.
\end{teorema}
O teste vem de $a_n=S_n-S_{n-1}\to S-S=0$ quando a série converge. A recíproca é \emph{falsa}:
$a_n\to0$ não garante nada (série harmônica). Por exemplo, $\sum\frac{n}{2n+1}$ diverge, pois
os termos tendem a $\frac12$.

\begin{teorema}[Teste da integral]{2.4}
\index{teste!da integral}Se $a_n=f(n)$, com $f$ contínua, positiva e decrescente para
$x\ge N$, então $\sum a_n$ e $\int_N^\infty f(x)\,dx$ convergem ou divergem juntas.
\end{teorema}

\figtikz{\begin{tikzpicture}
\begin{axis}[painel,width=8cm,height=4.6cm,xmin=0,xmax=7.4,ymin=0,ymax=1.15,xtick={1,2,3,4,5,6,7},ytick={1}]
\foreach \k in {1,...,6}{\edef\tmp{\noexpand\draw[fill=verde!45,draw=laranja] (axis cs:\k,0) rectangle (axis cs:{\k+1},{1/((\k+1)^2)});}\tmp}
\addplot[azul,thick,domain=0.93:7.3,samples=100]{1/x^2};
\end{axis}\end{tikzpicture}}{Teste da integral para $\sum\frac1{n^2}$: os retângulos de altura
$\frac1{(k+1)^2}$ cabem sob $y=\frac1{x^2}$, logo
$\sum_{n=2}^\infty\frac1{n^2}\le\int_1^\infty\frac{dx}{x^2}=1$.}

\begin{teorema}[Séries $p$]{2.5}
\index{série!$p$}A série $\displaystyle\sum_{n=1}^\infty\frac1{n^p}$ converge se $p>1$ e diverge
se $p\le1$.
\end{teorema}
De fato, $\int_1^\infty x^{-p}dx$ converge exatamente quando $p>1$ (capítulo~7 do Cálculo~I);
para $p\le0$ os termos nem tendem a zero.

\textbf{Estimativa do resto.} Nas condições do teste da integral, o erro de parar na soma
parcial $S_N$ é $R_N=S-S_N$, e
\[
\int_{N+1}^\infty f(x)\,dx\ \le\ R_N\ \le\ \int_N^\infty f(x)\,dx .
\]

\begin{exemplo}{2.4}{Quantos termos?}
Para $\sum\frac1{n^3}$, $R_N\le\int_N^\infty\frac{dx}{x^3}=\frac1{2N^2}$. Para garantir erro
menor que $0{,}001$, basta $2N^2>1000$, isto é, $N=23$. (O valor é $1{,}2020569\dots$, a
constante de Apéry.)
\end{exemplo}

% ---------------------------------------------------------------------
\section{Testes de comparação}

Quando os termos se parecem com os de uma série conhecida (geométrica ou $p$), comparamos.

\begin{teorema}[Teste da comparação]{2.6}
\index{teste!da comparação (séries)}Suponha $0\le a_n\le b_n$ para todo $n\ge N$. Se $\sum b_n$
converge, então $\sum a_n$ converge; se $\sum a_n$ diverge, então $\sum b_n$ diverge.
\end{teorema}

\begin{teorema}[Teste da comparação no limite]{2.7}
\index{teste!da comparação no limite}Sejam $a_n,b_n>0$ e $L=\lim\frac{a_n}{b_n}$. Se
$0<L<\infty$, as duas séries convergem ou divergem juntas. Se $L=0$ e $\sum b_n$ converge,
$\sum a_n$ converge; se $L=\infty$ e $\sum b_n$ diverge, $\sum a_n$ diverge.
\end{teorema}

\begin{exemplo}{2.5}{Comparando}
(a) $\sum\frac1{n^2+1}$: $\frac1{n^2+1}<\frac1{n^2}$, e $\sum\frac1{n^2}$ converge. Converge.

(b) $\sum\frac{\ln n}{n}$: para $n\ge3$, $\frac{\ln n}n>\frac1n$. Diverge.

(c) $\sum\frac1{2^n-1}$: a desigualdade vai no sentido errado para comparar com
$\frac1{2^n}$, mas $\frac{1/(2^n-1)}{1/2^n}\to1$. Pela comparação no limite, converge.

(d) $\sum\frac{3n^2+n}{\sqrt{n^7+2}}$: os termos se comportam como $\frac{3n^2}{n^{7/2}}=\frac3{n^{3/2}}$
(compare com $b_n=\frac1{n^{3/2}}$, $L=3$). Converge.
\end{exemplo}

\begin{nota}
Na comparação no limite, escolha $b_n$ guardando só os termos dominantes do numerador e do
denominador. A comparação direta exige a desigualdade no sentido certo; na dúvida, use o limite.
\end{nota}

% ---------------------------------------------------------------------
\section{Séries alternadas}

Uma \textbf{série alternada}\index{série!alternada} tem termos de sinais alternados:
$\sum(-1)^{n+1}b_n=b_1-b_2+b_3-\cdots$, com $b_n>0$.

\begin{teorema}[Teste da série alternada (Leibniz)]{2.8}
\index{teste!da série alternada}Se $b_{n+1}\le b_n$ (a partir de algum $n$) e $b_n\to0$, a série
$\sum(-1)^{n+1}b_n$ converge. Além disso, o erro de parar em $S_N$ satisfaz
\[
|R_N|=|S-S_N|\le b_{N+1}.
\]
\end{teorema}

\figtikz{\begin{tikzpicture}
\begin{axis}[painel,width=8cm,height=4.4cm,xmin=0,xmax=13,ymin=0.4,ymax=1.08,xtick={2,4,6,8,10,12},ytick={0.5,0.6931,1},yticklabels={$0{,}5$,$\ln2$,$1$},xlabel={$k$}]
\addplot[cinza,dashed,domain=0:13]{0.6931};
\addplot[laranja,thin] coordinates {(1,1)(2,0.5)(3,0.8333)(4,0.5833)(5,0.7833)(6,0.6167)(7,0.7595)(8,0.6345)(9,0.7456)(10,0.6456)(11,0.7365)(12,0.6532)};
\addplot[only marks,mark=*,mark size=1.4pt,azul] coordinates {(1,1)(2,0.5)(3,0.8333)(4,0.5833)(5,0.7833)(6,0.6167)(7,0.7595)(8,0.6345)(9,0.7456)(10,0.6456)(11,0.7365)(12,0.6532)};
\end{axis}\end{tikzpicture}}{Somas parciais da série harmônica alternada
$1-\frac12+\frac13-\cdots$: elas oscilam em torno de $\ln2$, com amplitude cada vez menor,
e a soma fica sempre entre duas somas parciais consecutivas.}

\begin{definicao}
$\sum a_n$ converge \textbf{absolutamente}\index{convergência absoluta} se $\sum|a_n|$ converge,
e \textbf{condicionalmente}\index{convergência condicional} se $\sum a_n$ converge mas
$\sum|a_n|$ diverge. Toda série absolutamente convergente é convergente.
\end{definicao}

\begin{exemplo}{2.6}{Absoluta ou condicional?}
(a) $\sum\frac{(-1)^{n+1}}n$ converge (Leibniz), mas $\sum\frac1n$ diverge: convergência
\emph{condicional}. Sua soma é $\ln2$.

(b) $\sum\frac{(-1)^{n+1}}{n^2}$ converge \emph{absolutamente}. Para erro menor que $0{,}01$,
basta $\frac1{(N+1)^2}<0{,}01$, isto é, $N=10$.

(c) $\sum\frac{\sen n}{n^2}$ não é alternada, mas $\big|\frac{\sen n}{n^2}\big|\le\frac1{n^2}$:
converge absolutamente.
\end{exemplo}

Séries condicionalmente convergentes são delicadas: reordenando seus termos é possível obter
qualquer soma desejada (teorema de Riemann). Com convergência absoluta, qualquer reordenação dá
a mesma soma.

% ---------------------------------------------------------------------
\section{Testes da razão e da raiz}

\begin{teorema}[Teste da razão]{2.9}
\index{teste!da razão}Seja $\rho=\lim_{n\to\infty}\Big|\dfrac{a_{n+1}}{a_n}\Big|$. Se $\rho<1$,
$\sum a_n$ converge absolutamente; se $\rho>1$ (ou $\rho=\infty$), diverge; se $\rho=1$, o teste
é inconclusivo.
\end{teorema}

\begin{teorema}[Teste da raiz]{2.10}
\index{teste!da raiz}Seja $\rho=\lim_{n\to\infty}\sqrt[n]{|a_n|}$. As mesmas conclusões do teste
da razão valem.
\end{teorema}

A ideia é que, para $n$ grande, a série se comporta como uma geométrica de razão $\rho$. O teste
da razão é o ideal para fatoriais e exponenciais; o da raiz, para potências $n$-ésimas.

\begin{exemplo}{2.7}{Razão e raiz}
(a) $\sum\frac{2^n}{n!}$: $\frac{a_{n+1}}{a_n}=\frac{2}{n+1}\to0$. Converge (para $e^2-1$,
como veremos no capítulo~3).

(b) $\sum\frac{n!}{n^n}$: $\frac{a_{n+1}}{a_n}=\frac{(n+1)\,n^n}{(n+1)^{n+1}}=\big(\frac n{n+1}\big)^n\to\frac1e<1$. Converge.

(c) $\sum\Big(\frac{2n+3}{3n+2}\Big)^n$: $\sqrt[n]{a_n}=\frac{2n+3}{3n+2}\to\frac23<1$. Converge.

(d) Para $\sum\frac1n$ e $\sum\frac1{n^2}$, $\rho=1$: um diverge e o outro converge, e o teste
não distingue.
\end{exemplo}

\begin{estrategia}{escolher um teste}
\begin{enumerate}[label=\arabic*.]
\item Os termos não tendem a zero? \emph{Diverge} (teste da divergência).
\item É geométrica, $p$ ou telescópica? Use o resultado conhecido.
\item Fração de polinômios ou raízes? \emph{Comparação} (no limite) com uma série $p$.
\item Tem $n!$ ou $c^n$? \emph{Razão}. Tem uma potência $(\cdots)^n$? \emph{Raiz}.
\item Sinais alternados? \emph{Leibniz}, e verifique se a convergência é absoluta.
\item $f(n)$ fácil de integrar (como $\frac1{n\ln n}$)? \emph{Integral}.
\end{enumerate}
\end{estrategia}

% ---------------------------------------------------------------------
\section*{Resumo do capítulo}\markright{Resumo do capítulo}
\addcontentsline{toc}{section}{Resumo do capítulo}
\begin{itemize}
\item Sequência converge se $a_n\to L$; monótona e limitada $\Rightarrow$ converge.
\item Série: limite das somas parciais. Geométrica: $\frac a{1-r}$ se $|r|<1$. Harmônica diverge.
\item Divergência: $a_n\not\to0\Rightarrow$ diverge (a recíproca é falsa).
\item Integral: $\sum f(n)$ e $\int f$ juntas; séries $p$ convergem sse $p>1$; resto entre duas integrais.
\item Comparação direta e no limite com séries geométricas e $p$.
\item Alternadas: Leibniz, erro $\le b_{N+1}$; convergência absoluta $\times$ condicional.
\item Razão e raiz: $\rho<1$ converge, $\rho>1$ diverge, $\rho=1$ nada se conclui.
\end{itemize}

% Gerado por ferramentas/gerar.py; não edite à mão.
\section*{Exercícios}\markright{Exercícios}
\addcontentsline{toc}{section}{Exercícios}
Os pontos antes de cada número indicam a dificuldade: \nivel{1} fácil, \nivel{2} médio, \nivel{3} difícil e \nivel{4} desafio. Os exercícios marcados com \estrela{} têm solução completa, com animação do gráfico, no repositório do livro: \url{https://github.com/AleffAzevedo/calculo-2-pt/blob/main/exercicios/cap2/README.md}. As respostas estão no fim do capítulo.

\begin{itemize}[leftmargin=5.2em,labelsep=0.6em,itemsep=4pt]
\item[\exnum{1}{2.1}] Calcule $\displaystyle\lim_{n\to\infty}\frac{2n^2+3n}{5n^2-1}$.
\item[\exnum{1}{2.2}] Calcule $\displaystyle\lim_{n\to\infty}n\sen\frac1n$.
\item[\exnum{1}{2.3}] Calcule $\displaystyle\lim_{n\to\infty}\Big(1+\frac2n\Big)^{3n}$.
\item[\exnum{1}{2.4}] Calcule $\displaystyle\lim_{n\to\infty}\big(\sqrt{n^2+n}-n\big)$.
\item[\exnum{1}{2.5}] A sequência $a_n=(-1)^n\dfrac n{n+1}$ converge?
\item[\exnum{1}{2.6}] Calcule $\displaystyle\lim_{n\to\infty}\frac{n^2}{2^n}$.
\item[\exnum{1}{2.7}] Calcule $\displaystyle\sum_{n=0}^\infty\Big(-\frac23\Big)^n$.
\item[\exnum{1}{2.8}] Calcule $\displaystyle\sum_{n=1}^\infty\frac{5\cdot3^n}{4^{n+1}}$.
\item[\exnum{1}{2.9}] Escreva $1{,}2\overline{3}=1{,}2333\ldots$ como fração usando uma série geométrica.
\item[\exnum{1}{2.10}] Calcule $\displaystyle\sum_{n=1}^\infty\frac1{(n+1)(n+2)}$ (telescópica).
\item[\exnum{1}{2.11}] Mostre que $\displaystyle\sum_{n=1}^\infty\ln\frac n{n+1}$ diverge calculando as somas parciais.
\item[\exnum{1}{2.12}] Use o teste da divergência em $\sum\frac{n^2}{n^2+1}$ e em $\sum\cos\frac1n$.
\item[\exnum{1}{2.13}] Classifique: $\sum\frac1{n^{1{,}01}}$, $\sum\frac1{\sqrt n}$, $\sum n^{-\pi}$.
\item[\exnum{1}{2.14}] Use o teste da integral para decidir sobre $\displaystyle\sum_{n=2}^\infty\frac1{n\ln n}$.
\item[\exnum{1}{2.15}] Use o teste da integral em $\sum n\,e^{-n^2}$.
\item[\exnum{1}{2.16}] Use comparação direta: $\sum\frac1{n^3+n}$ e $\sum\frac{2+\sen n}{n}$.
\item[\exnum{1}{2.17}] Use comparação no limite em $\displaystyle\sum\frac{n+1}{n^3+2}$.
\item[\exnum{1}{2.18}] Classifique $\displaystyle\sum\frac{(-1)^n}{\sqrt n}$ (absoluta, condicional ou divergente).
\item[\exnum{1}{2.19}] Use o teste da razão em $\displaystyle\sum\frac n{3^n}$.
\item[\exnum{1}{2.20}] Use o teste da razão em $\displaystyle\sum\frac{3^n}{n^2}$.
\item[\exnum{2}{2.21}] \estrela\ \textbf{Raízes encaixadas.} Seja $a_1=\sqrt2$ e $a_{n+1}=\sqrt{2+a_n}$, isto é, $\sqrt2,\ \sqrt{2+\sqrt2},\ \sqrt{2+\sqrt{2+\sqrt2}},\dots$
Mostre que a sequência converge e calcule o limite.
\item[\exnum{2}{2.22}] \estrela\ \textbf{A bola que quica.} Uma bola cai de $10$ m e, a cada quique, sobe a $\frac34$ da altura anterior. Qual a distância
total percorrida? E o tempo total, se a queda de uma altura $h$ leva $\sqrt{2h/g}$ ($g=9{,}8$)?
\item[\exnum{2}{2.23}] \estrela\ \textbf{A lentidão da série harmônica.} Use o teste da integral para mostrar que $\ln(n+1)\le H_n\le1+\ln n$, em que $H_n=1+\frac12+\dots+\frac1n$.
Quantos termos são necessários para $H_n$ passar de $10$?
\item[\exnum{2}{2.24}] \estrela\ \textbf{Telescópica com salto.} Calcule $\displaystyle\sum_{n=1}^\infty\frac1{n(n+2)}$.
\item[\exnum{2}{2.25}] \estrela\ \textbf{Estimando $\sum 1/n^2$.} Calcule $S_{10}$ para $\sum\frac1{n^2}$ e use as cotas do resto pelo teste da integral para obter um intervalo para a soma.
Compare com o valor exato $\frac{\pi^2}6$.
\item[\exnum{2}{2.26}] \estrela\ \textbf{Uma alternada que converge depressa.} A série $\displaystyle\sum_{n=1}^\infty\frac{(-1)^{n+1}}{n!}$ vale $1-\frac1e$. Quantos termos garantem erro menor que $10^{-4}$?
Compare a soma parcial com o valor exato.
\item[\exnum{2}{2.27}] Calcule $\displaystyle\sum_{n=1}^\infty\frac{2^n+3^n}{6^n}$.
\item[\exnum{2}{2.28}] Para quais $x$ a série $\sum_{n=0}^\infty(x-1)^n$ converge? Qual a soma?
\item[\exnum{2}{2.29}] Calcule $\displaystyle\lim_{n\to\infty}\sqrt[n]n$.
\item[\exnum{2}{2.30}] Mostre que $\dfrac{n!}{n^n}\le\dfrac1n$ e conclua o limite.
\item[\exnum{2}{2.31}] Calcule $\displaystyle\sum_{n=2}^\infty\frac1{n^2-1}$.
\item[\exnum{2}{2.32}] Decida: $\displaystyle\sum\frac{\arctg n}{n^2}$.
\item[\exnum{2}{2.33}] Decida: $\displaystyle\sum_{n=2}^\infty\frac1{n(\ln n)^2}$.
\item[\exnum{2}{2.34}] Decida: $\displaystyle\sum\frac{n^2+1}{n^4-n+3}$.
\item[\exnum{2}{2.35}] Decida: $\displaystyle\sum\sen\frac1n$.
\item[\exnum{2}{2.36}] Classifique $\displaystyle\sum\frac{(-1)^nn}{n^2+1}$.
\item[\exnum{2}{2.37}] Decida: $\displaystyle\sum\frac{(-1)^nn!}{10^n}$.
\item[\exnum{2}{2.38}] Decida: $\displaystyle\sum\frac{n^n}{n!}$.
\item[\exnum{2}{2.39}] Decida com o teste da raiz: $\displaystyle\sum\Big(\frac n{n+1}\Big)^{n^2}$.
\item[\exnum{2}{2.40}] Decida: $\displaystyle\sum\frac1{3^n-n}$.
\item[\exnum{3}{2.41}] \estrela\ \textbf{O floco de neve de Koch.} Partindo de um triângulo equilátero de área $A_0$, a cada etapa cada lado é dividido em três e, sobre o terço
do meio, constrói-se um novo triângulo equilátero para fora. Mostre que o perímetro tende a infinito, mas a área tende a $\frac85A_0$.
\item[\exnum{3}{2.42}] \estrela\ \textbf{Reordenar muda a soma.} Sabendo que $1-\frac12+\frac13-\frac14+\dots=\ln2$, mostre que a reordenação
$1+\frac13-\frac12+\frac15+\frac17-\frac14+\cdots$ (dois positivos, um negativo) tem soma $\frac32\ln2$.
\item[\exnum{3}{2.43}] \estrela\ \textbf{Empilhando blocos.} Blocos de comprimento $1$ são empilhados na borda de uma mesa. Colocando o centro de massa dos $k$ blocos de cima
sobre a borda do bloco $k+1$, o $k$-ésimo bloco (de cima para baixo) avança $\frac1{2k}$. Qual o balanço total com $n$ blocos?
Quantos blocos são precisos para o bloco do topo ficar inteiramente fora da mesa? E para um balanço de $3$?
\item[\exnum{3}{2.44}] \estrela\ \textbf{O conjunto de Cantor.} Do intervalo $[0,1]$ retira-se o terço do meio $(\frac13,\frac23)$; de cada um dos dois intervalos restantes,
retira-se o terço do meio; e assim por diante. Qual o comprimento total retirado?
\item[\exnum{3}{2.45}] \estrela\ \textbf{Somando $\sum n/2^n$.} Mostre que $\displaystyle\sum_{n=1}^\infty\frac n{2^n}$ converge e calcule a soma usando $S-\frac12S$.
Faça o mesmo para $\displaystyle\sum_{n=1}^\infty\frac{n^2}{2^n}$.
\item[\exnum{3}{2.46}] Decida: $\displaystyle\sum_{n=3}^\infty\frac1{n\ln n\,\ln(\ln n)}$.
\item[\exnum{3}{2.47}] Para quais $p$ a série $\displaystyle\sum_{n=2}^\infty\frac1{n(\ln n)^p}$ converge?
\item[\exnum{3}{2.48}] Decida: $\displaystyle\sum\frac{\ln n}{n^{1{,}5}}$.
\item[\exnum{3}{2.49}] Decida: $\displaystyle\sum\frac1{n^{1+1/n}}$.
\item[\exnum{3}{2.50}] Decida: $\sum\big(\sqrt{n+1}-\sqrt n\big)$ e $\sum\dfrac{\sqrt{n+1}-\sqrt n}{n}$.
\item[\exnum{3}{2.51}] Classifique $\displaystyle\sum_{n=2}^\infty\frac{(-1)^n}{\ln n}$.
\item[\exnum{3}{2.52}] Classifique $\displaystyle\sum(-1)^n\Big(1-\cos\frac1n\Big)$.
\item[\exnum{3}{2.53}] Decida: $\displaystyle\sum\frac{n^{100}}{1{,}01^n}$.
\item[\exnum{3}{2.54}] Decida: $\displaystyle\sum\frac{(n!)^2}{(2n)!}$.
\item[\exnum{3}{2.55}] Decida: $\displaystyle\sum\frac{1\cdot3\cdot5\cdots(2n-1)}{2\cdot4\cdot6\cdots(2n)}$. Por que o teste da razão não basta?
\item[\exnum{3}{2.56}] Seja $a_1=1$ e $a_{n+1}=\dfrac n{2n+1}a_n$. A série $\sum a_n$ converge?
\item[\exnum{3}{2.57}] Calcule $\displaystyle\sum_{n=1}^\infty\frac1{n(n+1)(n+2)}$.
\item[\exnum{3}{2.58}] Mostre que $\sum\frac{(-1)^n}{n^p}$ converge para todo $p>0$ e que a convergência é absoluta se, e só se, $p>1$.
\item[\exnum{3}{2.59}] Quantos termos de $\sum\frac{(-1)^{n+1}}{n^3}$ garantem erro menor que $0{,}0005$? Calcule a aproximação.
\item[\exnum{3}{2.60}] Seja $a_1=\tfrac12$ e $a_{n+1}=a_n(2-a_n)$. Mostre que a sequência é crescente e limitada por $1$ e calcule o limite.
\item[\exnum{4}{2.61}] Mostre que $\displaystyle\sum_{n=1}^k\frac1{n^2}<2-\frac1k$ usando $\frac1{n^2}<\frac1{n(n-1)}$ e conclua que $\sum\frac1{n^2}\le2$.
\item[\exnum{4}{2.62}] Seja $c_n=1+\frac12+\dots+\frac1n-\ln n$. Mostre que $\{c_n\}$ é decrescente e limitada inferiormente, logo converge (o limite é a constante de Euler $\gamma\approx0{,}5772$).
\item[\exnum{4}{2.63}] Usando $e=\sum_{n=0}^\infty\frac1{n!}$, prove que $e$ é irracional. (Suponha $e=\frac pq$ e multiplique por $q!$.)
\item[\exnum{4}{2.64}] (Condensação de Cauchy) Se $a_n\ge0$ é decrescente, mostre que $\sum a_n$ converge se, e só se, $\sum2^ka_{2^k}$ converge. Use isso para refazer as séries $p$.
\item[\exnum{4}{2.65}] Se $a_n>0$ e $\sum a_n$ converge, prove que $\sum a_n^2$ converge. Mostre com um exemplo que $\sum\sqrt{a_n}$ pode divergir.
\end{itemize}
\subsection*{Respostas}\markright{Respostas}
{\small\begin{multicols}{2}\raggedright
\begin{itemize}[leftmargin=2.7em,labelsep=0.5em,itemsep=1pt]
\item[\textbf{2.1}] $\tfrac25$.
\item[\textbf{2.2}] $1$ (com $t=\frac1n$, $\frac{\sen t}t\to1$).
\item[\textbf{2.3}] $e^6$.
\item[\textbf{2.4}] $\tfrac12$ (multiplique pelo conjugado).
\item[\textbf{2.5}] Não: os termos pares tendem a $1$ e os ímpares a $-1$.
\item[\textbf{2.6}] $0$ (L'Hôpital duas vezes em $\frac{x^2}{2^x}$).
\item[\textbf{2.7}] $\tfrac35$.
\item[\textbf{2.8}] $\tfrac{15}4$.
\item[\textbf{2.9}] $1{,}2+\tfrac{0{,}03}{1-0{,}1}=\tfrac{37}{30}$.
\item[\textbf{2.10}] $\tfrac12$.
\item[\textbf{2.11}] $S_k=-\ln(k+1)\to-\infty$.
\item[\textbf{2.12}] Ambas divergem (termos $\to1$).
\item[\textbf{2.13}] Converge; diverge; converge.
\item[\textbf{2.14}] $\int_2^\infty\tfrac{dx}{x\ln x}=\infty$: diverge.
\item[\textbf{2.15}] $\int_1^\infty xe^{-x^2}dx=\tfrac1{2e}$: converge.
\item[\textbf{2.16}] Converge ($<\tfrac1{n^3}$); diverge ($\ge\tfrac1n$).
\item[\textbf{2.17}] Com $\tfrac1{n^2}$, $L=1$: converge.
\item[\textbf{2.18}] Condicionalmente convergente.
\item[\textbf{2.19}] $\rho=\tfrac13$: converge.
\item[\textbf{2.20}] $\rho=3$: diverge.
\item[\textbf{2.21}] Crescente e limitada por $2$; o limite é $2$.
\item[\textbf{2.22}] $70$ m; cerca de $19{,}9$ s.
\item[\textbf{2.23}] $12\,367$ termos.
\item[\textbf{2.24}] $\tfrac34$.
\item[\textbf{2.25}] $S_{10}\approx1{,}5498$; $1{,}6407\le S\le1{,}6498$; $\frac{\pi^2}6\approx1{,}6449$.
\item[\textbf{2.26}] $N=7$; $S_7\approx0{,}632143$ contra $0{,}632121$.
\item[\textbf{2.27}] $\tfrac12+1=\tfrac32$.
\item[\textbf{2.28}] $0<x<2$; soma $\tfrac1{2-x}$.
\item[\textbf{2.29}] $1$ ($\ln\sqrt[n]n=\tfrac{\ln n}n\to0$).
\item[\textbf{2.30}] $\tfrac{n!}{n^n}=\tfrac1n\cdot\tfrac2n\cdots\tfrac nn\le\tfrac1n\to0$.
\item[\textbf{2.31}] $\tfrac12\big(\tfrac1{n-1}-\tfrac1{n+1}\big)$: soma $\tfrac34$.
\item[\textbf{2.32}] Converge ($\le\tfrac{\pi/2}{n^2}$).
\item[\textbf{2.33}] Converge ($\int_2^\infty=\tfrac1{\ln2}$).
\item[\textbf{2.34}] Converge (como $\tfrac1{n^2}$).
\item[\textbf{2.35}] Diverge (comparação no limite com $\tfrac1n$).
\item[\textbf{2.36}] Condicionalmente convergente.
\item[\textbf{2.37}] Diverge: os termos crescem sem limite.
\item[\textbf{2.38}] $\rho=e>1$: diverge.
\item[\textbf{2.39}] $\sqrt[n]{a_n}=\big(\tfrac n{n+1}\big)^n\to\tfrac1e<1$: converge.
\item[\textbf{2.40}] Converge (comparação no limite com $\tfrac1{3^n}$).
\item[\textbf{2.41}] Perímetro $3\ell(4/3)^n\to\infty$; área $A_0\big(1+\sum_{n\ge1}\frac{3\cdot4^{n-1}}{9^n}\big)=\frac85A_0$.
\item[\textbf{2.42}] $\tfrac32\ln2\approx1{,}0397$.
\item[\textbf{2.43}] $\frac12H_n$; $4$ blocos; $227$ blocos.
\item[\textbf{2.44}] $1$: o conjunto que sobra tem comprimento total zero.
\item[\textbf{2.45}] $2$ e $6$.
\item[\textbf{2.46}] Diverge (integral: $\ln\ln\ln x\to\infty$).
\item[\textbf{2.47}] $p>1$.
\item[\textbf{2.48}] Converge: $\ln n<n^{0{,}25}$ para $n$ grande, compare com $n^{-1{,}25}$.
\item[\textbf{2.49}] Diverge: $\frac{1/n^{1+1/n}}{1/n}=n^{-1/n}\to1$.
\item[\textbf{2.50}] A primeira diverge ($S_k=\sqrt{k+1}-1$); a segunda converge (termos $\sim\tfrac1{2n^{3/2}}$).
\item[\textbf{2.51}] Condicionalmente convergente ($\tfrac1{\ln n}>\tfrac1n$).
\item[\textbf{2.52}] Absolutamente convergente ($1-\cos\tfrac1n\sim\tfrac1{2n^2}$).
\item[\textbf{2.53}] $\rho=\tfrac1{1{,}01}<1$: converge (a exponencial vence).
\item[\textbf{2.54}] $\rho=\tfrac14$: converge.
\item[\textbf{2.55}] $\rho=1$; mas $a_n\ge\tfrac12\prod_{k=2}^n\tfrac{2k-2}{2k}=\tfrac1{2n}$: diverge por comparação com a harmônica.
\item[\textbf{2.56}] $\rho=\tfrac12$: converge.
\item[\textbf{2.57}] $\tfrac12\big(\tfrac1{n(n+1)}-\tfrac1{(n+1)(n+2)}\big)$: soma $\tfrac14$.
\item[\textbf{2.58}] Leibniz para $p>0$; $\sum\tfrac1{n^p}$ converge sse $p>1$.
\item[\textbf{2.59}] $\tfrac1{(N+1)^3}<0{,}0005$: $N=12$; $S_{12}\approx0{,}9013$ (valor $0{,}9015$).
\item[\textbf{2.60}] $1-a_{n+1}=(1-a_n)^2$; limite $1$.
\item[\textbf{2.61}] $S_k\le1+\sum_{n=2}^k\big(\tfrac1{n-1}-\tfrac1n\big)=2-\tfrac1k$.
\item[\textbf{2.62}] $c_n-c_{n+1}=\ln\tfrac{n+1}n-\tfrac1{n+1}>0$; e $c_n\ge\ln(n+1)-\ln n>0$ pelo teste da integral.
\item[\textbf{2.63}] $q!\,e-q!\sum_{n\le q}\tfrac1{n!}$ seria inteiro positivo menor que $\tfrac1{q+1}+\tfrac1{(q+1)^2}+\dots=\tfrac1q\le1$.
\item[\textbf{2.64}] Agrupe em blocos de $2^k$ termos; para $a_n=n^{-p}$, $\sum(2^{1-p})^k$ é geométrica.
\item[\textbf{2.65}] $a_n<1$ a partir de algum $n$, logo $a_n^2<a_n$; exemplo $a_n=\tfrac1{n^2}$.
\end{itemize}\end{multicols}}

